\documentclass[conference]{IEEEtran}
\usepackage{cite}
\usepackage{amsmath,amssymb,amsfonts}
\usepackage{algorithmic}
\usepackage{graphicx}
\usepackage{textcomp}
\usepackage{xcolor}
\usepackage{booktabs}
\usepackage{tabularx}
\usepackage{balance}
\usepackage{xurl}

% =====================================================================
%  REVISED DRAFT - 2026-09-22 (With Advanced LLM Metrics)
%  Framing: descriptive longitudinal case study + exploratory triangulation.
% =====================================================================

\begin{document}
\raggedbottom

\title{Do We Still Need to Be Agile? What Two Cohorts of Student Projects Reveal About Planning Discipline in the GenAI-Era}

\author{\IEEEauthorblockN{Anderson Martins Gomes}
\IEEEauthorblockA{Department of Computer Science \\
State University of Cear\'a (UECE) \\
Fortaleza, Cear\'a, Brazil \\
anderson.gomes@aluno.uece.br}
}

\maketitle

\begin{abstract}
In 2001, the Agile Manifesto declared war on upfront planning. A
quarter-century later, Generative AI has inverted the economics of
software production: generating code is cheap, while architectural 
alignment and human coordination remain scarce and expensive. We tracked
14 team-semesters (9 teams in semester 2025.2, 5 teams in 2026.1) of
software-engineering project courses to examine how planning discipline 
affects late-stage integration. By leveraging an LLM to evaluate $T_1$ 
commit messages and parsing file-provenance to separate deferred development 
from destructive rework, we found that Agile's de-emphasis on planning 
was taken to an extreme: 55\% of the first cohort completely omitted 
substantive early architectural alignment. While repository metadata 
spiked late at $T_3$, LLM-derived qualitative transcript analysis revealed 
that human coordination friction remained chronically high (8/10) from the 
very beginning. Critically, teams with low upfront planning quality were 
heavily penalized with catastrophic late-stage rework during integration 
(up to 100,000+ lines), whereas high-planning teams contained rework to 
nominal levels despite producing massive volumes of deferred code. We 
present Spec-Driven Development (SDD) not merely as a proposal, but as a 
data-backed necessity for the AI era: when generating code is trivial, 
upfront machine-readable specifications act as the primary shield against 
destructive rework.
\end{abstract}

\begin{IEEEkeywords}
Generative AI, Agile Software Development, Planning Debt, Code Churn,
Empirical Software Engineering, Software Engineering Education.
\end{IEEEkeywords}

\section{Introduction}

The Agile Manifesto revolutionized software engineering by prioritizing
working software, individual interactions, and responding to change
over comprehensive documentation and strict upfront planning. For over
two decades, this empirical paradigm was the gold standard for
navigating requirements uncertainty, and for good reason: in a world
where typing and boilerplate were the bottlenecks, minimizing upfront
documentation minimized waste.

That world is gone. Generative AI assistants now produce thousands of
lines of code in seconds, shifting the scarce resources of software
production directly to architectural alignment, integration, and
cognitive coordination. Under this new cost equation, the Agile bet ---
that planning shrinks as understanding grows --- deserves to be tested
rather than assumed. 

If deferring design in an AI-assisted environment merely relocates effort, 
we hypothesize a collision between students' initial hubris regarding AI-driven 
productivity and the harsh reality of late-stage architectural integration. 
This paper tests that intuition by triangulating repository metadata, 
multi-checkpoint student surveys, and qualitative meeting transcripts across 
two student cohorts. We structure our investigation around a ``perception versus 
reality'' arc, asking three questions:

\begin{itemize}
    \item \textbf{RQ1 (Descriptive):} How do student developers utilize Generative AI, and how does this adoption shape their longitudinal perceptions of software engineering roles and project expectations?
    \item \textbf{RQ2 (Descriptive):} How does the temporal density of repository activity contrast with the qualitative typification of human coordination friction across the project lifecycle?
    \item \textbf{RQ3 (Primary):} How does the quality of upfront architectural artifacts (evaluated at baseline $T_1$) associate with late-stage destructive rework and independent, blind-scored project outcomes?
\end{itemize}

This paper makes three primary contributions to the empirical understanding of software engineering processes in the GenAI era. First, we provide a mixed-methods characterization of the dissonance between developer expectations and integration realities. Utilizing longitudinal survey data and qualitative sentiment tracking, we map a distinct trajectory from initial technological hubris---where students presume AI will seamlessly automate implementation and replace coding roles---to late-stage disillusionment driven by chronic architectural and coordination friction (Section~\ref{sec:results}). 

Second, our primary empirical contribution isolates the mechanical and evaluative costs of deferred planning. To overcome the unreliability of student commit messages, we employ Large Language Models (LLMs) to conduct an artifact state analysis directly on the repository baseline at $T_1$. By cross-referencing this upfront architectural clarity with file provenance (separating deferred new development from destructive rework churn) and double-blind evaluator scorecards, we provide concrete evidence that planning quality governs integration stability. Teams exhibiting poor upfront artifact quality incurred massive rework churn and suffered severe penalties in peer-reviewed project evolution scores, whereas high-planning teams effectively contained AI-generated code volumes without destabilizing the core architecture (Section~\ref{sec:results}).

Finally, we synthesize these empirical findings into a data-backed methodological proposition. Rather than presenting Spec-Driven Development (SDD) as a mere theoretical alternative to Agile empiricism, we demonstrate it as a structural necessity. Because Generative AI drives the cost of code generation toward zero while severely compounding contextual integration friction, establishing explicit, machine-readable specifications upfront serves as a mandatory countermeasure to safely absorb AI-generated throughput and prevent late-stage project collapse (Section~\ref{sec:discussion}).

\section{Background and Related Work}

\subsection{The Agile value hierarchy and coordination risk}
The Agile Manifesto~\cite{agile_manifesto_2001} established a value hierarchy that 
prioritized working software and human adaptability over comprehensive documentation. 
In pre-AI environments, where typing code and boilerplate served as the physical 
bottleneck of production, minimizing upfront documentation minimized waste. 
However, critics have long maintained that empirical, documentation-light processes 
do not eliminate coordination costs; they merely relocate them. Brooks's law regarding 
the communication overhead of late software projects~\cite{brooks_1975_mythical_man_month}, and 
DeMarco and Lister's findings on social and defensive barriers in 
\emph{Peopleware}~\cite{demarco_lister_1987_peopleware}, serve as canonical counterpoints. 
When system structure is communicated primarily through tacit team agreement rather 
than explicit artifacts, the project remains highly vulnerable to cognitive overload 
during late-stage integration.

\subsection{Generative AI in SE: From syntax to systems}
Recent systematic reviews and empirical studies on Generative AI in software engineering 
demonstrate a stark dichotomy in AI capabilities~\cite{nguyenduc2026systematic,ieee2025transforming}. 
While Large Language Models (LLMs) excel at localized code generation, 
increasing development speed significantly, their effectiveness in high-level 
architectural decisions remains fundamentally limited~\cite{ieee2025transforming}. 
The current literature highlights a severe ``contextual gap'': LLMs possess superhuman 
speed for generating logic, but lack persistent understanding of cross-file context 
and problem scoping~\cite{cambridge2026contextual}. Consequently, GenAI drastically 
accelerates localized implementation, effectively shifting the software production 
bottleneck away from syntax generation and directly onto human cognitive coordination 
and system-wide alignment.

\subsection{The resurgence of upfront planning: BDUF and SDD}
Historically, the software industry rejected ``Big Design Up Front'' (BDUF) because 
changing rigid, comprehensive architectural documents mid-project was economically 
disastrous when paired with slow, manual coding phases. However, Generative AI has 
inverted this economic equation. If AI can translate a design into working code in 
seconds, the cost of iterating on the implementation approaches zero, while the cost 
of unguided architectural drift rises exponentially. 

This shift motivates Spec-Driven Development (SDD) --- an approach where structured 
specifications are used as the primary source of truth for AI-assisted 
software~\cite{thoughtworks202Xspec}. In SDD, deterministic, machine-readable 
specifications govern probabilistic AI code generation, moving the engineering focus 
from ``code is the source of truth'' to ``intent is the source of truth''~\cite{github2025spec}. 
By strictly isolating architectural intent from execution, SDD represents a modern, 
high-speed incarnation of BDUF designed to enforce the system boundaries that 
autonomous tools otherwise degrade.

\subsection{Technical debt and integration rework}
Cunningham's debt metaphor~\cite{cunningham_1992_technical_debt} and Fowler's practitioner 
treatment~\cite{fowler2009technicaldebt} provide the language of principal and 
interest that we borrow for \emph{Planning Debt}: effort deferred during the design 
phase, which must be repaid with severe interest during late-stage integration. 
The technical-debt measurement literature~\cite{kruchten2012technicaldebt} cautions 
that blunt volume-based proxies often conflate productive agile iteration with 
destructive rework. This distinction is critical in the GenAI era, where AI-generated 
code is empirically shown to dramatically accelerate technical debt and double code 
churn~\cite{codebridge2026hidden}. This reality directly shaped our methodology; 
rather than relying on total lines changed, our pipeline analyzes file provenance 
to specifically isolate the destructive rework churn caused by architectural collapse 
at the project deadline.

\section{Methodology}
\label{sec:method}

To investigate how planning discipline and coordination evolve in AI-assisted software engineering projects, we designed an empirical study combining longitudinal repository mining, standardized instructor evaluations, multi-checkpoint student surveys, and qualitative meeting transcripts across two academic offerings.

\subsection{Context, Course Design, and "Rules of the Game"}
The study was conducted within a project-based undergraduate software engineering course at the State University of Ceará (UECE), spanning two semesters: 2025.2 ($N = 9$ analyzable teams) and 2026.1 ($N = 5$ teams), yielding a total inferential dataset of $N = 14$ team-semesters ($\approx 70$ participating students). 

Students formed teams of approximately five members through spontaneous self-allocation, with remaining students assigned to vacant slots to ensure balanced team sizes. Teams enjoyed complete autonomy regarding their internal organization, process models, and project management methodologies (e.g., choosing freely between Agile variants, Scrum, Kanban, or unguided approaches). However, course guidelines enforced two structural constraints: 
\begin{enumerate}
    \item Teams were required to develop an operational software project featuring an integrated Generative AI component as a core functional requirement. Specifically, projects were expected to implement advanced architectural patterns such as Retrieval-Augmented Generation (RAG)---a technique that dynamically grounds Large Language Models (LLMs) in external knowledge bases to mitigate hallucination \cite{lewis2020retrieval} --- or LLM-driven data analytics pipelines.
    \item Students were strongly encouraged to utilize generative AI assistants throughout all phases of software development, from requirements specification to implementation and testing.
\end{enumerate}

The project lifecycle was anchored by three synchronized temporal checkpoints ($T_1$, $T_2$, and $T_3$) paced identically across both cohorts:
\begin{itemize}
    \item \textbf{$T_1$ (Baseline):} Administered on the first day of class prior to any formal instruction on generative AI processes, capturing baseline perceptions, initial project scope, and early planning commits.
    \item \textbf{$T_2$ (Mid-Semester):} Conducted following structured instructional modules covering generative AI mechanics, RAG architectures, and software process models, accompanied by the second wave of survey instruments and status presentations.
    \item \textbf{$T_3$ (Final Integration):} Positioned one week prior to the final project delivery and presentation, capturing the peak pressure of system integration and deadline convergence.
\end{itemize}

\subsection{Project Scope and Real-World Complexity}
To ensure ecological validity and induce realistic coordination friction, the course guidelines explicitly prohibited trivial algorithmic exercises or "toy" applications. As part of the "rules of the game," teams were mandated to identify and solve real-world community, educational, or business problems. This constraint forced students to grapple with high domain complexity, authentic user requirements, and the nontrivial challenge of embedding Generative AI into functional architectures.

The resulting software systems spanned diverse domains, ranging from healthcare and financial management to academic research and ecotourism. Table~\ref{tab:projects} details the 14 projects analyzed in this study. The breadth and complexity of these applications validate that the observed late-stage integration challenges and chronic coordination friction were driven by authentic software engineering workloads rather than artificially simplistic academic tasks. 

\begin{table*}[t]
\caption{Analyzed Student Projects by Cohort, Demonstrating Real-World Problem Scope and Generative AI Integration}
\label{tab:projects}
\centering
\begin{tabular}{@{}lp{3cm}p{12.5cm}@{}}
\toprule
\textbf{Semester} & \textbf{Project Name} & \textbf{System Description \& AI Integration Scope} \\
\midrule
2026.1 & FoodGuard & Web platform utilizing nutritional data and medical history to identify food allergens and prevent adverse reactions. \\
2026.1 & RevisAI & Educational system that generates intelligent flashcards and organizes customized weekly study review cycles. \\
2026.1 & SmartFlowIA & SaaS WhatsApp chatbot focusing on intelligent onboarding, orientation, and welcoming for university freshmen. \\
2026.1 & Interview Prep IA & AI-driven system that analyzes resumes against job vacancies, simulates interviews, and provides personalized candidate feedback. \\
2026.1 & EnemBot & Domain-specific chatbot designed to tutor and assist students with questions regarding the Brazilian National High School Exam. \\
\midrule
2025.2 & UECEncontra & Intelligent virtual assistant and campus navigation system designed to help students and professors locate university resources. \\
2025.2 & EstudaAI & Study management platform that processes multimodal files (audio, image, PDF) to generate lecture summaries, reminders, and track grades. \\
2025.2 & CareerPath IA & Resume analysis tool that uses GenAI to match user profiles with market demands and suggest targeted skill development. \\
2025.2 & Buliçoso & Healthcare application that utilizes LLMs to simplify complex medical leaflets and manage medication reminders for the elderly. \\
2025.2 & EcoRoteiro & Ecotourism platform that generates personalized travel itineraries based on user profiles, physical conditions, and specific interests. \\
2025.2 & Research Flow & Academic research assistant that searches articles, generates intelligent summaries, and formats texts to academic standards (e.g., ABNT, APA). \\
2025.2 & ExamForge & Assessment simulator that ingests past exams and course materials to generate new tests matching specific professor grading patterns. \\
2025.2 & SoFiA & Financial management web platform utilizing AI to help university students dynamically organize their personal finances. \\
2025.2 & Synapse & Requirements engineering tool that summarizes technical client meetings to automatically generate functional requirements, user stories, and backlogs. \\
\bottomrule
\end{tabular}
\end{table*}

\subsection{Data Sources and Collection Framework}
Our data collection framework triangulates four distinct evidentiary channels to capture both technical artifacts and human behavioral dynamics.

\textbf{1. Git Repositories:} Complete historical commit logs, branch structures, line change counts, and file-level provenance were extracted from all team repositories across the entire semester lifecycle. To mitigate the threat of unobserved out-of-band collaboration, students were explicitly instructed to centralize all project artifacts---including requirements, architectural specifications, documentation, and source code---directly within their Git repositories. This directive ensured the repository served as a comprehensive, single source of truth for the team's entire production effort.

\textbf{2. Independent Evaluator Scorecards:} At each checkpoint ($T_1$, $T_2$, $T_3$), teams delivered formal status presentations. These were assessed by a panel of independent supervisors (4 evaluators in the 2025.2 cohort, and 3 in the 2026.1 cohort). To prevent groupthink and ensure scoring validity, an independent blind-scoring protocol was enforced: following the presentation and open feedback, each supervisor submitted their scores via a private electronic form. Evaluators remained blind to their peers' scores throughout the semester. The final project score was calculated as an unweighted average of these independent assessments. Projects were assessed across four structured criteria on explicit bounded scales: \emph{Project Evolution} (0 to 4.0 points), \emph{Scope and Applicability} (0 to 2.0 points), \emph{Technical Complexity} (0 to 2.0 points), and \emph{Engagement and Participation} (0 to 2.0 points).

\textbf{3. Multi-Checkpoint Student Surveys:} Students completed identical survey instruments at $T_1$, $T_2$, and $T_3$ (totaling six instrument cycles). These surveys systematically captured data across five core constructs: Demographics, Professional Experience, GenAI Adoption, Perceptions of AI Impact on SE Roles, and Project Expectations. To ensure replicability, the core dimensions and sample variables tracked across these checkpoints are detailed in Table~\ref{tab:survey_instrument}.

\begin{table*}[t]
\caption{Structure of the Multi-Checkpoint Survey Instrument}
\label{tab:survey_instrument}
\centering
\begin{tabularx}{\textwidth}{@{}>{\raggedright\arraybackslash}p{3.5cm}>{\raggedright\arraybackslash}X>{\raggedright\arraybackslash}p{2.5cm}@{}}
\toprule
\textbf{Construct} & \textbf{Evaluated Variables \& Inquiries} & \textbf{Data Type} \\
\midrule
\textbf{Demographics \& Academic Background} & Age, Gender, Current academic semester, Previous degrees, High school educational background. & Categorical \\
\addlinespace
\textbf{Professional SE Experience} & Years of software development experience, Current employment status, Preferred engineering role (e.g., Backend, QA, ML), Largest past team size. & Categorical \\
\addlinespace
\textbf{Generative AI Adoption \& Tooling} & Frequency of AI usage, Specific tools adopted (e.g., ChatGPT, Gemini, GitHub Copilot, Claude), Primary use cases (e.g., Code generation, Debugging, Documentation). & Categorical \& Multiple Choice \\
\addlinespace
\textbf{Perceived AI Impact on SE Roles} & Perceived risk of extinction or severe disruption for specific roles: Backend, Frontend, QA, Project Manager, Product Manager, Scrum Master. & 5-point Likert \\
\addlinespace
\textbf{Expectations \& Motivation} & Confidence levels prior to project initiation, Motivation to learn new techniques, Expected 5-year impact of AI on the industry. & 5-point Likert \& Open-ended \\
\bottomrule
\end{tabularx}
\end{table*}

\textbf{4. Qualitative Meeting Transcripts:} To capture coordination dynamics and evaluative feedback, audio recordings were collected from the three formal status presentation checkpoints. This dataset comprises the full audio records of the student presentations alongside the subsequent deliberative discussions among the supervisors in a private chat group. Crucially, these supervisory discussions were free-form qualitative reflections, completely decoupled from the numeric scorecasting process. To make the transcription data computationally feasible for LLM analysis, the raw text was segmented into analytical chunks, yielding 100 distinct fragments for the 2025.2 cohort (35 at $T_1$, 37 at $T_2$, and 28 at $T_3$).

\textbf{Data Anonymization and Privacy:} To ensure strict ethical compliance and protect participant confidentiality, all raw data streams were subjected to a rigorous anonymization pipeline prior to analysis. We employed a hybrid sanitization approach, combining deterministic heuristics (e.g., regex-based masking of identifiers) with LLM-assisted redaction to scrub unstructured text across commit messages, survey free-text responses, and audio transcripts. Consequently, all Personally Identifiable Information (PII) was completely removed, ensuring that no references to individual students, evaluators, or specific demographic markers remain in the finalized analytical dataset.

\subsection{Advanced Metrics Operationalization}
\label{sec:advanced-metrics}

Sections~\ref{sec:method}A--D established \emph{what} was collected. This
subsection formalizes \emph{how} raw Git history, evaluator scorecards, and
survey or transcript text become the nine metrics (M1--M9) used for
RQ1--RQ3. Each metric below states, in order: the raw input, the exact
transformation applied to it, and the resulting value. No new data are
collected and no new LLM calls are issued beyond the fixed-rubric scoring
already described in this section and in Section~\ref{sec:advanced-metrics}'s
provenance notes.

% PROVENANCE: All nine scripts below live under paper_v8/scripts/metrics/,
% write CSV artifacts to paper_v8/data/, and read only already-persisted,
% frozen Phase 1/2 lake and analysis Parquet files or the frozen
% paper_v4/advanced_metrics/outputs/*.csv signals produced by
% 12_paper_signals_extractor.py. No script reruns Phase 1/2 or issues new
% LLM calls; all LLM scoring referenced here was already computed upstream.

\subsubsection{RQ1: Generative-AI adoption and role-perception metrics}

RQ1 is operationalized entirely from self-reported survey text and Likert
items, at cohort~$\times$~cut grain (survey responses are anonymized and
carry no team identifier). We do not treat any Git-derived signal as
evidence of AI-tool usage: as detailed under M3 (Section~\ref{sec:advanced-metrics},
RQ2), our repository data can identify \emph{which teammate} authored a
commit and \emph{how concentrated} authorship was before a deadline, but
cannot determine whether a generative-AI tool assisted that commit. This
absence of a verified behavioral AI-usage signal is treated as an explicit
construct-validity limitation rather than papered over with a proxy.

\textbf{M1 -- Self-Reported AI-Dependency Trajectory.} Each student's
free-text answer to six survey prompts (perceived AI benefit, five-year
career-impact expectation, autonomy/tool-dependency balance, career
expectations, anticipated project challenges, and feelings about the
project) is independently rated by a large language model on an ordinal
$[0,4]$ ``AI-dependency'' scale, one rating per individual response, using a
fixed rubric prompt at zero decoding temperature. For cohort-cut $(s,c)$,
each question family's per-response ratings are averaged, and the composite
trajectory value is the unweighted mean across the six family means:
\[
\mathrm{AIDependency}(s,c) = \frac{1}{|Q|} \sum_{q \in Q} \overline{\mathrm{ai\_dependency\_score}}_{q,s,c}.
\]
% PROVENANCE: source data/analysis/textual_cut_signals.parquet (Semestre x
% temporal_marker grain, no ID_Equipe; produced by
% 04_nlp_qualitative_miner.py's aggregation stage, itself built from
% per-response LLM scores using pipeline_prompts.py's STUDENT_NLP_PROMPT,
% model gpt-4o-mini, temperature 0). Computed by
% paper_v8/scripts/metrics/m1_ai_dependency_trajectory.py, writing
% paper_v8/data/m1_ai_dependency_trajectory.csv. This script only averages
% already-scored, already-aggregated per-family means; it issues no new LLM
% calls.

\textbf{M2 -- Perceived Role-Disruption Risk.} For each surveyed SE role
(backend, frontend, QA, project manager, product manager, and Scrum Master)
and cohort-cut $(s,c)$, we summarize the raw 5-point Likert item (``this role
will be extinct or severely affected by generative AI tools,'' 1~=~strongly
disagree, 5~=~strongly agree) exactly as submitted by students, with no LLM
or other transformation involved:
\[
\mathrm{RoleRisk}(r,s,c) = \mathrm{SummarizeDistribution}\big(\{\ell_{r,s,c}\}\big),
\]
where $\mathrm{SummarizeDistribution}(\cdot)$ returns mean, median, IQR, and
missingness over the raw item values $\ell_{r,s,c}$.
% PROVENANCE: source data/lake/student_responses.parquet (Semestre x
% temporal_marker grain, no ID_Equipe; validated via
% phase2_contracts.load_phase2_inputs, contract "student_responses").
% Computed by paper_v8/scripts/metrics/m2_role_disruption_risk.py, writing
% paper_v8/data/m2_role_disruption_risk.csv. New metric, zero LLM calls
% (raw structured Likert responses).

\subsubsection{RQ2: Repository temporal density vs. qualitative coordination friction}

\textbf{M3 -- Pre-Deadline Author-Concentration Density.} For team $i$ in
semester $s$, we take every commit in the 72-hour window immediately
preceding that team's $T_3$ checkpoint and group it by the committing
teammate's identifier. Let $C(a,i,s)$ be teammate $a$'s commit count in that
window and $C(i,s)$ the window's total commit count; we report the median
per-teammate commit share and the Gini coefficient of the per-teammate
commit-count distribution:
\[
\begin{aligned}
\mathrm{AuthorShare}(i,s)
  &= \underset{a}{\mathrm{median}}
     \left(\frac{C(a,i,s)}{C(i,s)}\right), \\
\mathrm{AuthorGini}(i,s)
  &= \mathrm{Gini}\bigl(\{C(a,i,s)\}_a\bigr).
\end{aligned}
\]
We emphasize what this metric can and cannot support: it identifies how
unevenly commit volume was distributed across identified human teammates
immediately before the deadline; it does \emph{not} identify whether any
individual commit involved generative-AI assistance, since no tool-identity
signal (author name/email pattern, commit-message marker, or similar) is
present in our collected data. We therefore report M3 as an
integration-pressure proxy for RQ2, not as AI-adoption evidence for RQ1.
% PROVENANCE: source data/analysis/integration_friction_metrics.parquet
% (team-semester grain; the underlying per-teammate commit counts, shares,
% and Gini coefficient are computed directly from the raw ID_Autor_Local
% field in data/lake/git_commits.parquet, restricted to a fixed 72-hour
% pre-T3 window; no AI/human classification is applied anywhere in this
% computation). Re-exposed, not recomputed, by
% paper_v8/scripts/metrics/m3_author_concentration_density.py, writing
% paper_v8/data/m3_author_concentration_density.csv.

\textbf{M4 -- Repository Activity Density.} For every commit, we take the
per-file line-addition and line-deletion counts reported by a standard Git
diff and sum them into a single per-commit churn value (binary file changes,
for which line-level diffs are undefined, are counted separately as binary
events rather than contributing to churn). For semester $s$ (and pooled
across both semesters) and cut $c \in \{T_1,T_2,T_3\}$, we summarize the
cross-team distribution of two aggregates: total churn
$\mathrm{cc\_total}$ (summed over all commits in the cut) and commit count
$\mathrm{cc\_commit\_n}$. Here $m$ denotes either aggregate:
\[
\mathrm{AD}(s,c,m)
  = \mathcal{S}\bigl(\{m_{i,s,c}:i\in\mathrm{teams}(s)\}\bigr),
\]
where $\mathcal{S}(\cdot)$ returns mean, median, IQR, and missingness. No
file-path or extension filtering is applied to these totals.
% PROVENANCE: source data/analysis/code_churn_metrics.parquet (team-semester
% grain; computed from per-commit lines_added/lines_deleted sums in
% data/lake/git_commits.parquet, with binary events tracked separately per
% pipeline_config.py's excluded-from-line-churn policy). Computed by
% paper_v8/scripts/metrics/m4_repo_activity_density.py, writing
% paper_v8/data/m4_repo_activity_density.csv.
Empirically, mean $\mathrm{cc\_total}$ escalates by roughly two orders of
magnitude from $T_1$ to $T_3$ in both cohorts (Section~\ref{sec:results}),
confirming that repository-visible activity concentrates sharply at the
deadline.

\textbf{M5 -- Qualitative Coordination-Friction Trajectory.} For each cut
$c \in \{T_1,T_2,T_3\}$, we pool the full transcript text of every recorded
status-presentation session for that cut, across all teams and both
semesters, into one text block. Each of the three resulting blocks is rated
once by a large language model on a $[1,10]$ coordination-friction scale
(considering blockers, misunderstandings, handoff problems, merge/integration
conflicts, and repeated coordination overhead), using a fixed rubric prompt
at zero decoding temperature, judged only from the supplied text.
% PROVENANCE: source paper_v4/advanced_metrics/outputs/cohort_temporal_friction.csv
% (Semestre x temporal_marker grain, Semestre pooled as "all"; one LLM call
% per cut, three calls total, model gpt-4o-mini, temperature 0, prompt
% "cohort_coordination_friction" in paper_v4/advanced_metrics/llm_prompts.py;
% produced by 12_paper_signals_extractor.py). Re-exposed, not re-scored, by
% paper_v8/scripts/metrics/m5_coordination_friction_trajectory.py, writing
% paper_v8/data/m5_coordination_friction_trajectory.csv.
The contrast between M4 and M5 is the core empirical basis for RQ2: unlike
$\mathrm{ActivityDensity}$, which is near zero at $T_1$ and spikes at $T_3$,
$\mathrm{CoordinationFriction}$ is already at its ceiling-adjacent value at
$T_1$ and remains flat through $T_3$ (Section~\ref{sec:results}). Repository
metadata is therefore a lagging indicator of a coordination cost that, per
the qualitative record, was present from day one.

\subsubsection{RQ3: Upfront planning quality, destructive rework, and evaluator outcomes (primary)}

\textbf{M6 -- $T_1$ Planning-Quality Score.} For team $i$ in semester $s$, we
take the subject line of every commit made before that team's $T_1$
checkpoint, concatenate them in chronological order into a single text
block, and rate that block once with a large language model on a $\{1,\dots,10\}$
scale of clarity, specificity, and coherence, using a fixed rubric prompt at
zero decoding temperature, judged only from the supplied text.
$\mathrm{PlanningScore}(i,s)$ is undefined (missing) when a team made zero
commits before $T_1$, since there is no text to rate.
% PROVENANCE: source paper_v4/advanced_metrics/outputs/team_level_signals.csv
% (team-semester grain; one LLM call per team-semester, model gpt-4o-mini,
% temperature 0, prompt "planning" in
% paper_v4/advanced_metrics/llm_prompts.py; produced by
% 12_paper_signals_extractor.py). Re-exposed by
% paper_v8/scripts/metrics/m6_t1_planning_quality.py, writing
% paper_v8/data/m6_t1_planning_quality.csv.

\textbf{M7 -- $T_1$ Planning-Omission Rate.} For cohort $s$, the omission
rate is the share of teams for which $\mathrm{PlanningScore}$ (M6) is
undefined, i.e., teams with no commits at all before $T_1$:
\[
\mathrm{OmissionRate}(s)
=\frac{\left|\left\{i\in\mathrm{teams}(s):P(i,s)=\mathrm{NaN}\right\}\right|}
       {|\mathrm{teams}(s)|}.
\]
% PROVENANCE: derived deterministically (no LLM calls) from
% paper_v4/advanced_metrics/outputs/team_level_signals.csv by
% paper_v8/scripts/metrics/m7_planning_omission_rate.py, writing
% paper_v8/data/m7_planning_omission_rate.csv.
Applying this to the collected commit records gives an omission rate of
$55.6\%$ ($5/9$) for 2025.2 and $0\%$ ($0/5$) for 2026.1. The former is the
source of the abstract's ``55\% of the first cohort'' claim.

\textbf{M8 -- Rework Ratio.} For each team, we exclude boilerplate paths from
the recorded file-level changes. The filters cover dependency directories
(\path{node_modules}, \path{venv}, and \path{.venv}), lockfiles, JSON data,
images, icons, caches, source maps, and minified artifacts. For every
remaining file, we record the earliest checkpoint ($T_1$, $T_2$, or $T_3$) at
which it appears in that team's history. Restricting to file-change records
observed at $T_3$, we then sum lines added plus lines deleted separately for
files whose earliest checkpoint was $T_1$ or $T_2$ (destructive rework of
pre-existing work) versus files first seen at $T_3$ itself (deferred new
development):
\[
\begin{aligned}
\mathrm{TotalChurn}_{T_3}(i,s)
  &= \mathrm{ReworkChurn}_{T_3}(i,s) \\
  &\quad + \mathrm{DeferredChurn}_{T_3}(i,s),
\end{aligned}
\]
\[
\mathrm{ReworkRatio}_{T_3}(i,s) = \frac{\mathrm{ReworkChurn}_{T_3}(i,s)}{\mathrm{TotalChurn}_{T_3}(i,s)}.
\]
% PROVENANCE: rework_churn_t3 / deferred_churn_t3 columns from
% paper_v4/advanced_metrics/outputs/team_level_signals.csv (produced by
% 12_paper_signals_extractor.py's file-provenance step, applied to
% data/lake/git_files.parquet, cross-checked against
% data/processed/clean_repos/). Derived deterministically by
% paper_v8/scripts/metrics/m8_rework_severity_ratio.py, writing
% paper_v8/data/m8_rework_severity_ratio.csv.
$\mathrm{ReworkRatio}_{T_3}$ isolates \emph{what fraction} of late-cycle churn
is destructive rather than raw volume, making teams of very different scale
comparable.

\textbf{M9 -- Planning Quality vs. Rework/Outcome Association.} We join, at
team-semester grain, the $T_1$ Planning-Quality Score (M6), the Rework
Severity Ratio and rework-churn volume (M8), and four independent evaluator
outcomes at $T_3$ (project progress, scope/applicability, technical
complexity, and engagement/participation, each the arithmetic mean of the
numeric scores independently submitted per team-cut by every evaluator on
the panel, on the bounded scales already described in
Section~\ref{sec:method}). Given $n \leq 14$, we report two descriptive
views: (a) Spearman rank correlation $\rho\big(\mathrm{PlanningScore}, Y\big)$
for each outcome $Y$; and (b) a median-split contrast of
$\mathrm{PlanningScore}$ (teams with an undefined score, per M7, are excluded
from the split and reported separately, not imputed), summarizing each $Y$
per group with the same distribution-summary procedure used throughout this
subsection.
% PROVENANCE: joins paper_v4/advanced_metrics/outputs/team_level_signals.csv
% with data/analysis/cross_evidence/datasets/evaluator_outcome_metrics.parquet
% (team-semester grain, itself an arithmetic-mean aggregation of raw
% per-evaluator form submissions in data/lake/evaluator_team_cuts.parquet).
% Computed by paper_v8/scripts/metrics/m9_planning_vs_rework_association.py,
% writing paper_v8/data/m9_planning_vs_rework_association.csv and
% paper_v8/data/m9_planning_vs_rework_association_group_contrast.csv. Uses
% scipy.stats.spearmanr (scipy==1.11.4, already pinned in requirements.txt).
Consistent with this project's own exploratory-evidence stance (Threats to
Validity), no correlation in (a) is presented as more than descriptive
association at $n=9$ scored teams; the median-split contrast in (b) is the
quantitative basis for the abstract's claim that low-planning teams incurred
disproportionate $T_3$ rework (up to $100{,}000+$ lines for a single team)
while high-planning teams contained rework to comparatively nominal levels
(Section~\ref{sec:results}).

Table~\ref{tab:metrics_summary} summarizes M1--M9. RQ1 relies exclusively on
self-reported survey signals (M1, M2), since no
verified behavioral AI-usage signal exists in our repository data. RQ2
contrasts a repository-derived integration-pressure proxy (M3) and
quantitative activity density (M4) against qualitative coordination friction
(M5) on a shared $T_1\text{--}T_3$ axis. RQ3 forms a single evidentiary chain
at team-semester grain: an upfront planning signal (M6), its cohort-level
omission rate (M7), a late-stage destructive-rework signal (M8), and their
joint association with evaluator outcomes (M9). Concrete values, figures, and
the RQ3 group-contrast findings are reported in Section~\ref{sec:results}.

\begin{table*}[t]
\caption{Summary of Metrics M1--M9, Grouped by Research Question}
\label{tab:metrics_summary}
\centering
\begin{tabularx}{\textwidth}{@{}>{\raggedright\arraybackslash}p{0.7cm}>{\raggedright\arraybackslash}p{4.4cm}>{\raggedright\arraybackslash}p{2.6cm}>{\raggedright\arraybackslash}X@{}}
\toprule
\textbf{ID} & \textbf{Metric} & \textbf{Grain} & \textbf{Source Artifact} \\
\midrule
\multicolumn{4}{@{}l}{\textit{RQ1 -- GenAI adoption and role-perception}} \\
\midrule
M1 & Self-Reported AI-Dependency Trajectory & Cohort $\times$ cut & \path{textual_cut_signals.parquet} \\
M2 & Perceived Role-Disruption Risk & Cohort $\times$ cut & \path{student_responses.parquet} \\
\midrule
\multicolumn{4}{@{}l}{\textit{RQ2 -- Repository density vs. coordination friction}} \\
\midrule
M3 & Pre-Deadline Author-Concentration Density & Team-semester & \path{integration_friction_metrics.parquet} \\
M4 & Repository Activity Density & Semester $\times$ cut & \path{code_churn_metrics.parquet} \\
M5 & Coordination-Friction Trajectory & Cohort $\times$ cut & \path{cohort_temporal_friction.csv} \\
\midrule
\multicolumn{4}{@{}l}{\textit{RQ3 -- Planning quality, destructive rework, and outcomes (primary)}} \\
\midrule
M6 & $T_1$ Planning-Quality Score & Team-semester & \path{team_level_signals.csv} \\
M7 & $T_1$ Planning-Omission Rate & Cohort & \path{team_level_signals.csv} \\
M8 & Rework Severity Ratio & Team-semester & \path{team_level_signals.csv} \\
M9 & Planning vs. Rework/Outcome Association & Team-semester & \path{team_level_signals.csv} + \path{evaluator_outcome_metrics.parquet} \\
\bottomrule
\end{tabularx}
\end{table*}

\section{Results}
\label{sec:results}

[to be written...]

\section{Discussion}
\label{sec:discussion}

[to be written...]

\section{Threats to Validity}
[to be rewritten...]
\textbf{Construct.} We used GPT-4o mini as a qualitative coder for planning
quality and coordination friction. Although the prompts were structured and the
temperature was set to zero, LLM evaluations retain interpretative variance.

\textbf{Internal.} Observational design: no intervention, no randomization. 
GenAI exposure was self-reported and treated as the underlying context of the 
cohorts rather than a manipulated variable.

\textbf{External.} $N=14$ team-semesters from two semesters of one course; 
student greenfield micro-team context. Generalization to industrial settings 
requires further study.

\section{Conclusion}
[to be written...]

\bibliographystyle{IEEEtran}
\bibliography{references}

\end{document}